\clearpage
\SetRegion{第四章 全球开源技术全景洞察}{rust}{network}
\phantomsection\pdfbookmark[1]{本章导览：从功能分工进入技术生态}{ch4-guide}
\kicker{04 / GLOBAL OPEN SOURCE LANDSCAPE}
\PageTitle{从功能分工进入技术生态}{先看完整全景，再追问资源如何被组织和复用。}
技术生态由不同能力共同支撑：任务入口连接工具与模型，模型运行依托数据、计算和工程系统，开放规范使多种实现能够继续协作。本章以六个观察入口呈现这些分工，在完整名单中定位具体对象，再选择有证据支持的问题深入分析。

\noindent\includegraphics[width=\linewidth]{figures/new/chapter4_navigation.pdf}
\vspace{2mm}

4.1、4.2分别从项目协作和模型获取信号进入分析；4.3—4.6保留完整技术全景。全景的分区与位置用于功能导航，不表示真实依赖、成熟程度或市场规模。每个条目链接到对应仓库或模型页面。

\noindent\includegraphics[width=\linewidth]{figures/new/chapter4_dimensions.pdf}
\src{与本章直接相关的四个发展维度，用于提出问题；不将项目收录数量转换为目标达成评分。}
\ReportNote{主统计期：2025-07-01至2026-06-30。4.1采用期后固定名单回看主期；4.2的2026-10-02模型快照及HF夏季观察分别标注。其他全景沿用固定对象与分类，具体时点、待核状态随图说明。项目Logo与账户头像用于辨识对象，不构成托管、许可或认证证明。}
